\section{Posterior predictive model criticism: derivations}
\label{sec:annex-p15-posterior-predictive-model-criticism}

\begin{annexitem}{annex:p15.ppc-joint}{thm:p15-ppc-joint}{Constructing the paired reference law}
\begin{proof}
Fix the observed data and the specified version of every conditional kernel.
For rectangles define
\[
 Q_y(C\times B)=\int_C R_{A(y)}(\theta,B)\Pi_y(d\theta).
\]
The probability-kernel product construction gives a probability measure
on the product sigma algebra; its total mass is the integral of one.
Taking \(C=\Theta\) gives the stated data marginal. Joint measurability
of \(D\) makes the comparison event measurable. Integrating its indicator
first over the replication kernel and then over the posterior gives exactly
\(\int q_D(y,\theta)\Pi_y(d\theta)\), by Tonelli's theorem.
If \(D=T\), the event depends only on the replicated data and the fixed
threshold \(T(y)\), proving the statistic case.

The importance of sharing the parameter can be demonstrated without an
asymptotic argument. Let the data space contain one point, let the posterior
be uniform on \(\{0,1\}\), and put \(D(y,\theta)=\theta\).
The shared-parameter comparison is always true, so its probability is one.
If independent posterior parameters are used in the two discrepancies,
the inequality fails when the observed-side parameter is one and the
replicated-side parameter is zero, so its probability is only \(3/4\).
Thus the two constructions can differ even with exactly sampled posteriors.
\end{proof}
\end{annexitem}

\begin{annexitem}{annex:p15.ppc-simulation}{thm:p15-ppc-simulation}{Sampling error and conditional averaging}
\begin{proof}
Under the fixed-data reference law, each \(I_j\) is Bernoulli with
success probability \(p_B(y)\). Independence of pairs gives the mean
and variance of their average. Also
\(E(I_j\mid\theta_j,y)=q_D(y,\theta_j)\), so the tower identity
gives the mean of \(\widetilde p\). Conditional variance decomposition
gives
\[
 \operatorname{Var}(I_j\mid y)=
 E[q_D(y,\theta_j)(1-q_D(y,\theta_j))\mid y]
 +\operatorname{Var}(q_D(y,\theta_j)\mid y).
\]
The first term is nonnegative; divide by \(J\) to obtain the variance
comparison for independent pairs. Chebyshev's inequality and
\(p(1-p)\leq1/4\) prove the probability bound.
For stationary indicators, expand
\(\operatorname{Var}(\sum_jI_j)=\sum_{i,j}\operatorname{Cov}(I_i,I_j)\).
There are \(J\) diagonal terms and \(2(J-h)\) terms at lag \(h\),
which gives the displayed covariance formula. A process not stationary
with the target marginal need not have mean \(p_B(y)\); no assertion
about its bias or convergence follows from that covariance identity.
\end{proof}
\end{annexitem}

\begin{annexitem}{annex:p15.ppc-pivot}{thm:p15-ppc-pivot}{Pivotal calibration and conditional contraction}
\begin{proof}
If every conditional statistic law is \(F_a\), integrating that law
against any posterior leaves it unchanged. For a continuous distribution
function \(F\), the probability integral transform sends a variable
with that distribution to a uniform variable. One proof uses its quantiles:
for \(0<u<1\), continuity and the generalized inverse identify a
threshold with cumulative probability \(u\); any flat intervals have
probability zero. Hence the probability that \(F(T)\leq u\) is \(u\).
Reflecting about \(1/2\) gives the same statement for \(1-F(T)\).

For the second assertion, let \(F_\theta\) be the continuous
distribution function of \(D(Y,\theta)\) under \(P_\theta\), and put
\(U=1-F_\vartheta(D(Y,\vartheta))\). Conditional on \(\vartheta\), \(U\)
is uniform; therefore its unconditional mean is \(1/2\) and variance
is \(1/12\). The absence of atoms makes an inclusive tail equal to
\(1-F_\theta\) at almost every observed discrepancy. Conditional
independence of the replicate yields
\[
 p_B(Y)=E[U\mid Y].
\]
Taking expectations proves the mean identity. Finally,
\(\operatorname{Var}(U)=E[\operatorname{Var}(U\mid Y)]
+\operatorname{Var}(E[U\mid Y])\) proves the variance bound.
The joint prior law was used in identifying the posterior as the conditional
law of \(\vartheta\); this argument is not a fixed-\(\theta\) calibration
claim. If there are ties, the inclusive-tail conditional expectation is
instead \((1+\Pr(D(Y',\vartheta)=D(Y,\vartheta)))/2\) after averaging
the conditionally exchangeable pair, so continuity cannot simply be dropped.
\end{proof}
\end{annexitem}

\begin{annexitem}{annex:p15.ppc-normal-mean}{thm:p15-ppc-normal-mean}{The normal mean check and its sampling distribution}
\begin{proof}
The likelihood factor depending on \(\theta\) is
\(\exp\{-(\theta-\bar y)^2/(2s^2)\}\). Multiplication by the
normal prior and completion of the square give posterior precision
\(1/s^2+1/\tau^2\), variance \(v\), and mean \(v\bar y/s^2=w\bar y\).
A replicated sample mean equals a posterior draw of \(\theta\) plus
an independent \(N(0,s^2)\) error. Its law is therefore
\(N(w\bar y,v+s^2)\), proving the upper-tail formula.

Under the proper prior model, \(\bar Y\sim N(0,\tau^2+s^2)\).
Consequently the argument of \(\Phi\) in that formula is centered
normal with variance
\[
 \frac{(1-w)^2(\tau^2+s^2)}{s^2+v}
 =\frac{s^2}{s^2+2\tau^2}<1.
\]
If \(1-\Phi(Z)\) were uniform, applying the inverse normal distribution
function would make \(Z\) standard normal, contradicting that variance.
With the flat prior the normalized likelihood is exactly
\(N(\bar y,s^2)\), so the predictive sample mean is
\(N(\bar y,2s^2)\). Its probability of exceeding its own center is
\(1/2\). This expression does not involve any other aspect of \(y\).
\end{proof}
\end{annexitem}

\begin{annexitem}{annex:p15.ppc-linear}{thm:p15-ppc-linear}{Orthogonal residual decomposition}
\begin{proof}
Full column rank makes \(B^TB\) positive definite. The matrix \(P\)
is symmetric, idempotent and has range equal to the column space of
\(B\). For \(\widehat\beta=(B^TB)^{-1}B^Ty\), orthogonality gives
\[
 \|y-B\beta\|^2
 =\|(I_n-P)y\|^2
  +(\beta-\widehat\beta)^TB^TB(\beta-\widehat\beta).
\]
The nonnegative second term vanishes at \(\widehat\beta\); this proves
the minimum formula. The flat-prior posterior density is proportional to
the exponential of minus half that quadratic term, hence is a proper
\(N_k(\widehat\beta,(B^TB)^{-1})\) density. Under it the quadratic
term has a \(\chi^2_k\) law and mean \(k\), proving the average formula.

Generate \(Y'=B\beta+\varepsilon'\), with
\(\varepsilon'\sim N_n(0,I_n)\) independent of the posterior draw.
Then \(r(Y')=\|(I_n-P)\varepsilon'\|^2\sim\chi^2_{n-k}\),
because an orthonormal basis diagonalizing the projection has \(n-k\)
residual coordinates and the normal law is invariant under orthogonal maps.
This distribution is independent of \(\beta\), so it is also the
predictive law. Since \(D_{\rm avg}(Y')=r(Y')+k\), both reduced
statistics give the same tail.

The realized replicated discrepancy is
\(D(Y',\beta)=\|\varepsilon'\|^2=U\sim\chi^2_n\).
The observed discrepancy is \(r(y)+V\), where the posterior quadratic
term \(V\sim\chi^2_k\) depends only on \(\beta\). Thus \(U,V\)
are independent, giving the claimed difference law.
For a standard normal variable, integrating its density by parts gives
\(E Z^2=1\), \(E Z^4=3\), and hence
\(\operatorname{Var}(Z^2)=2\). Summing independent squares gives the
stated means and variances, including the sum of the two variances for
the difference.

Both chi-square laws with positive degrees of freedom have positive
density on \((0,\infty)\). Independence therefore implies
\(\Pr(V>U)>0\), for example by taking \(U<1\) and \(V>2\).
At \(r=0\) the difference tail is strictly below one, while the
\(\chi^2_{n-k}\) tail is one. By continuity this strict ordering also
holds for sufficiently small positive \(r\). Thus an assertion that the
realized tail is always larger is false. Equality of means and a comparison
of variances do not establish stochastic or upper-tail dominance.
\end{proof}
\end{annexitem}

\begin{annexitem}{annex:p15.ppc-poisson}{thm:p15-ppc-poisson}{The total and the allocation contain different information}
\begin{proof}
For fixed data the product Poisson likelihood is proportional to
\(\lambda^S e^{-n\lambda}\). Multiplying by the gamma prior yields
the asserted posterior. Conditional on \(\lambda\), the replicate total
is Poisson with mean \(n\lambda\). Thus its predictive mass is
\[
 \frac{n^r(b+n)^{a+S}}{r!\Gamma(a+S)}
 \int_0^\infty \lambda^{a+S+r-1}e^{-(b+2n)\lambda}\,d\lambda.
\]
The substitution \(u=(b+2n)\lambda\) and the defining gamma integral
give the displayed formula. Its dependence on the observed vector is
only through \(S\).

For nonnegative integers \(x_i\) summing to \(S\), divide the joint
Poisson mass by the Poisson mass of its total:
\[
 \Pr(Y'=x\mid S'=S,\lambda)
 =\frac{e^{-n\lambda}\lambda^S/\prod_i x_i!}
        {e^{-n\lambda}(n\lambda)^S/S!}
 =\frac{S!}{\prod_i x_i!}n^{-S}.
\]
The parameter cancels; this is the claimed multinomial law, also after
any posterior averaging. With \(n=4,S=20\), direct expansion gives
\(T(x)=\sum_i x_i^2/5-20\). Since
\(\sum_i x_i^2\leq(\sum_i x_i)^2=400\), its maximum is 60.
Equality requires all cross-products \(x_ix_j\), \(i\ne j\), to
vanish, hence all 20 counts to occupy one of the four cells. Each such
allocation has mass \(4^{-20}\); the four are disjoint. The upper
tail at 60 is consequently \(4^{1-20}\). At the equal allocation the
nonnegative squared-residual sum is zero, so its inclusive upper tail is one.
\end{proof}
\end{annexitem}

\begin{annexitem}{annex:p15.ppc-mortality-posterior}{thm:p15-ppc-mortality-posterior}{Normalization on an inequality region}
\begin{proof}
The unnormalized density is nonnegative and measurable. Its integral over
\(C\) is bounded by the integral over the positive orthant:
\[
 Z\leq\prod_{t=35}^{64}\int_0^\infty
               \theta_t^{y_t}e^{-N_t\theta_t}\,d\theta_t
   =\prod_{t=35}^{64}\frac{\Gamma(y_t+1)}{N_t^{y_t+1}}<\infty.
\]
For positivity, choose a strictly positive, strictly increasing and strictly
convex sequence, for instance \(\theta_t=(t-34)^2\). All the finitely
many strict inequalities persist on a sufficiently small open neighborhood.
The density is continuous and positive there, so the integral over that
neighborhood is positive. This proves \(Z>0\). The product gamma
density on the orthant differs from the integrand only by a positive
constant; conditioning it on \(C\), which has positive probability,
therefore gives the normalized density claimed.

If an additional upper bound of one is imposed, a sufficiently small
positive multiple of the same strictly convex sequence supplies an interior
point of the bounded region. The same argument proves posterior positivity
and finiteness there as well. These are two explicitly specified regions,
not an assertion that their normalized posteriors are identical.
Finally, the sampling model at a fixed rate vector factors over the ages.
Holding that vector and the exposures fixed therefore gives independent
Poisson replication with the displayed means.
\end{proof}
\end{annexitem}

\begin{annexitem}{annex:p15.ppc-mortality-fit}{thm:p15-ppc-mortality-fit}{A positive-increment representation and its optimality conditions}
\begin{proof}
Index the rates by \(j=t-35\), and multiply them by 1000. The
coefficient \(\eta_0\) is the first scaled rate, \(\eta_1\) is
its first forward difference, and each later coefficient is the successive
increase in the forward difference. Telescoping the first differences
gives exactly the displayed positive-part formula. Nonnegative first
difference and nonnegative increments are equivalent to monotonicity and
discrete convexity. These observations prove existence and uniqueness
of the representation. Equivalently, the map from its thirty coefficients
to the thirty rates is an invertible triangular linear map. Multiplying
each rate by its positive exposure preserves invertibility.

For \(y>0\) and \(\mu>0\), the second derivatives of
\(\mu-y\log\mu\) and \((y-\mu)^2/\mu\) are
\(y/\mu^2\) and \(2y^2/\mu^3\), respectively. Thus both sums
are strictly convex in the mean vector and, by the invertible linear
map, in the coefficients. Each scalar summand is bounded below and tends
to infinity as \(\mu\) tends to either zero or infinity. It follows
that every nonempty sublevel set of either sum confines every mean to a
compact interval within \((0,\infty)\). Intersecting with the closed
linear inequalities defining the rate region yields a compact set.
There are feasible strictly positive rate vectors of finite objective
value. Continuity therefore gives attainment; strict convexity gives
uniqueness. Although the first coefficient must be positive, the sublevel
set is bounded away from its excluded zero boundary by the first positive
count's contribution.

At a positive coefficient, both signs of sufficiently small coordinate
perturbations are feasible, so differentiability at a minimum forces its
derivative to vanish. At a zero coefficient, an increasing perturbation
is feasible and forces a nonnegative derivative. Conversely, assume
these conditions at a feasible \(\eta\). For every other feasible
\(\xi\), the supporting-plane inequality for a differentiable convex
function gives
\[
 F(\xi)\geq F(\eta)+\sum_j
       \frac{\partial F}{\partial\eta_j}(\xi_j-\eta_j)
       \geq F(\eta).
\]
The positive-coordinate terms are zero; a zero-coordinate term is
nonnegative because both the derivative and \(\xi_j\) are nonnegative.
This proves sufficiency. Ordinary differentiation of each scalar summand,
followed by the chain rule for its linear mean, gives the two gradient
formulas. The proof uses the strictly positive counts in this data set;
zero counts need a separate boundary and uniqueness analysis.
\end{proof}
\end{annexitem}

\begin{annexitem}{annex:p15.ppc-mixture}{thm:p15-ppc-mixture}{Augmentation, deviance and label invariance}
\begin{proof}
The Dirichlet normalizer follows by integrating over the simplex.
For two coordinates the integral is
\(B(\kappa_1,\kappa_2)=\Gamma(\kappa_1)\Gamma(\kappa_2)/
\Gamma(\kappa_1+\kappa_2)\). For the induction step, put
\(x_1=t\) and \(x_i=(1-t)u_i\) for \(i\geq2\).
The Jacobian in the remaining \(m-2\) free coordinates is
\((1-t)^{m-2}\). The integral factors into the \((m-1)\)-coordinate
simplex integral and \(B(\kappa_1,\sum_{i\geq2}\kappa_i)\),
giving \(\prod_i\Gamma(\kappa_i)/\Gamma(\kappa_+)\).
The one-coordinate case is a point mass. Increasing \(\kappa_i\)
by one in this normalizing integral gives
\(E X_i=\kappa_i/\kappa_+\), justifying the stated center.

The complete multinomial probability for class and cell is
\(w_hu_{ha}v_{hb}z_{hc}\). Conditional on the observed cell, Bayes'
rule normalizes these quantities over \(h\), giving the latent-cell
allocation probabilities in the body. For augmented counts the factors
depending on parameters are
\[
 \prod_h w_h^{n_{h+++}}
 \prod_{h,a}u_{ha}^{n_{ha++}}
 \prod_{h,b}v_{hb}^{n_{h+b+}}
 \prod_{h,c}z_{hc}^{n_{h++c}},
\]
where a plus denotes summation over that index. Multiplication by the
independent Dirichlet densities adds these exponents to their respective
prior parameters minus one. This proves all conditional posterior laws,
including conditional independence of the blocks. All parameters remain
strictly positive, so the laws are proper. The observed-data posterior is
proper as well: the multinomial likelihood is a positive bounded function
on the interior of the probability simplexes and the prior is proper.

Write \(q_{abc}=y_{abc}/N\). Apart from the same multinomial coefficient,
the saturated likelihood is \(\prod q_{abc}^{y_{abc}}\) and the model
likelihood is \(\prod p_{abc}^{y_{abc}}\). Twice their log ratio is
exactly the stated deviance. For positive \(q_i\), the inequality
\(-\log x\geq1-x\) gives
\[
 \sum_{q_i>0}q_i\log(q_i/p_i)
 \geq\sum_{q_i>0}(q_i-p_i)
 =1-\sum_{q_i>0}p_i\geq0.
\]
If a required \(p_i\) vanishes the discrepancy is infinite, consistent
with nonnegativity. Permuting class indices changes no sum
\(\sum_h w_hu_{ha}v_{hb}z_{hc}\), proving invariance.
Generating \(N\) fresh independent class-and-cell draws and summing over
their class coordinates yields the multinomial law with these cell
probabilities. This generation resamples class membership; fixing the
original augmented counts would define a different conditional experiment.
\end{proof}
\end{annexitem}

\begin{annexitem}{annex:p15.ppc-normal-maximum}{thm:p15-ppc-normal-maximum}{Prior and posterior maximum distributions}
\begin{proof}
The identity
\(\sum_i(y_i-\theta)^2=\sum_i(y_i-\bar y)^2+n(\theta-\bar y)^2\)
shows that the likelihood times the bounded uniform prior is proportional
to \(\mathbf1_{[-L,L]}(\theta)e^{-n(\theta-\bar y)^2/2}\).
Its integral is finite and positive, proving the truncated normal formula.
Given \(\theta\), the probability that one replicate lies in \([-t,t]\)
is \(\Phi(t-\theta)-\Phi(-t-\theta)\). Independence gives its
\(n\)-th power for the event that the absolute maximum is at most \(t\).
Averaging over the posterior or the prior gives the two displayed laws.

For fixed \(t\geq0\), let
\(g_t(\theta)=\Phi(t-\theta)-\Phi(-t-\theta)\). Since
\(0\leq G_t=g_t^n\leq g_t\), Tonelli gives
\[
 \int_{\mathbb R}G_t(\theta)\,d\theta
 \leq\int_{-t}^t\int_{\mathbb R}\phi(x-\theta)\,d\theta\,dx=2t,
\]
where \(\phi\) is the standard normal density. Thus
\(F_{\rm prior}(t)\leq t/L\), proving its limit. Let \(Z\) have
the untruncated \(N(\bar y,1/n)\) law. For every bounded measurable
\(h\),
\[
 E[h(Z)\mid |Z|\leq L]
 =\frac{E[h(Z)\mathbf1\{|Z|\leq L\}]}{\Pr(|Z|\leq L)}
 \longrightarrow E[h(Z)]
\]
by dominated convergence and a denominator tending to one. Apply this
to \(h=G_t\). It proves the posterior predictive limit and identifies
the untruncated normal as a limit or approximation, not the exact posterior
for a bounded prior.
\end{proof}
\end{annexitem}
